\begin{center}\begin{tikzpicture}[font=\small,>=stealth,line join=round]
\draw[gray!40,rounded corners=3pt] (0,0) rectangle (14.2,8.3);\node at (7.1,7.700000000000001) {曲线可以靠近y轴，却没有x=0这个点};
\begin{scope}[shift={(1.6454545454545455,4.897058823529411)},x=2.1818181818181817cm,y=0.9019607843137255cm]\draw[->,gray] (-0.25,0)--(2.5,0) node[right]{$x$};\draw[->,gray] (0,-3.6)--(0,1.5) node[above]{$y$};\draw[blue!75!black,thick,solid,domain=0.125:2.2,samples=180] plot(\x,{ln(\x)/ln(2)});\fill[orange!85!black] (0.5,-1) circle(2.2pt);\fill[orange!85!black] (0.25,-2) circle(2.2pt);\fill[orange!85!black] (0.125,-3) circle(2.2pt);\draw[gray,dashed] (.125,-3)--(.125,0);\draw[gray!35,densely dotted] (1,0)--(1,-3.6);\node[below=3pt] at (1,-3.6){$1$};\draw[gray!35,densely dotted] (2,0)--(2,-3.6);\node[below=3pt] at (2,-3.6){$2$};\end{scope}\node[anchor=west,align=left,text width=5.0cm] at (8.4,4.0) {$x=\frac12:\ y=-1$\\[9pt]$x=\frac14:\ y=-2$\\[9pt]$x=\frac18:\ y=-3$\\[9pt]每个横坐标都大于0\\[9pt]永远不能令$x=0$};
\node[text width=12.6cm,align=center] at (7.1,.65) {靠近的是x=0这条竖直线，并非经过y轴上的某个点；定义域没有0。};
\end{tikzpicture}\end{center}
